\section{AI Critique}

Đối với sự tác động của AI trên nhiều lĩnh vực nói chung và lĩnh vực kiểm thử phần mềm nói riêng, ta thấy rằng hệ thống AI dường như vẫn còn nhiều sai lầm trong quá trình nhận diện yêu cầu và trả lời nội dung liên quan. Cụ thể hơn là sau khi tiếp nhận yêu cầu và đưa ra câu trả lời liên quan đến các chủ đề cần sự minh bạch về thông tin như thống kê , mã lỗi , v.v.. thì hệ thống AI còn nhiều sai sót ở việc đưa ra những số liệu không xác định được, hoặc những thông tin không có căn cứ rõ ràng. Điều này có thể dẫn đến việc người dùng bị hiểu lầm hoặc nhận được thông tin sai lệch, đặc biệt là trong các lĩnh vực đòi hỏi sự chính xác cao như y tế, tài chính, hay pháp luật. Lí do lớn nhất cũng như dễ nhận dạng nhất là hệ thống AI luôn tiếp nhận vô số dữ liệu từ nhiều nguồn khác nhau, và không phải tất cả dữ liệu đó đều được kiểm chứng hoặc có độ tin cậy cao. Do đó, việc AI đưa ra những thông tin không chính xác có thể là do nó đã học từ những dữ liệu không đáng tin cậy hoặc đã bị lỗi trong quá trình xử lý thông tin. Điều này nhấn mạnh tầm quan trọng của việc kiểm tra và xác minh thông tin mà AI cung cấp, cũng như việc cải thiện chất lượng dữ liệu mà AI sử dụng để học hỏi.
Từ những sai sót đã nêu, đối với bài tập này em thấy rằng để làm việc một cách hiệu quả với AI thì người dùng cần phải có khả năng đánh giá và kiểm tra thông tin mà AI cung cấp một cách cẩn thận. Đối với những thông tin liên quan đến tính xác thực cao mình cần phải đối chiếu với các nguồn tin cậy khác hoặc sử dụng các công cụ kiểm tra thông tin để đảm bảo rằng những gì AI cung cấp là phù hợp chứ không phải mang tính ảo giác hay thiên vị.